% !TEX root =  ../master.tex

\section*{Kurzfassung (Abstract)}

Diese Studienarbeit dokumentiert die methodische und technische Portierung interaktiver, Python-basierter Lehr-Notebooks der Vorlesung \enquote{Theoretische Informatik III} in das TypeScript-Ökosystem. Ziel war es, die didaktischen und architektonischen Vorteile statischer Typisierung für die Vermittlung formaler Konzepte nutzbar zu machen. Um die referenzbasierte Gleichheitsprüfung von JavaScript zu überwinden, wurde die performante Datenstruktur \texttt{RecursiveSet} entwickelt, die mathematische Wertsemantik und effizientes Hashing für verschachtelte Objekte über ein \textit{Structure of Arrays} (SoA) Speicherlayout garantiert. Darüber hinaus analysiert eine Case Study den Paradigmenwechsel vom Parsergenerator \emph{PLY} zur deklarativen, EBNF-nahen Syntaxanalyse mit \emph{Lezer}. Die Ergebnisse zeigen, dass strengere Typisierung das Verständnis für Compilerbau-Konzepte fördert, wenngleich dies eine höhere Code-Verbosität bedingt. Die Arbeit schließt mit einer Reflexion über den iterativen, KI-gestützten Übersetzungsprozess und bestätigt, dass eine erfolgreiche Migration stets eine strukturelle Neuausrichtung erfordert.

\vspace{1em} % Hier ist der saubere Abstand! (Kannst du auch auf 0.5cm o.ä. ändern)

\noindent This student research project documents the methodological and technical migration of interactive, Python-based Jupyter Notebooks for the course \enquote{Theoretical Computer Science III} into the TypeScript ecosystem. The objective was to leverage the didactic and architectural benefits of static typing for teaching formal concepts. To overcome JavaScript's reference-based equality evaluation, the high-performance data structure \texttt{RecursiveSet} was developed, guaranteeing mathematical value semantics and efficient hashing for nested objects via a Structure of Arrays (SoA) memory layout. Furthermore, a case study analyzes the paradigm shift from the \emph{PLY} parser generator to declarative, EBNF-like syntax analysis using \emph{Lezer}. The results demonstrate that stricter typing enhances the understanding of compiler construction concepts, despite requiring increased code verbosity. The paper concludes with a reflection on the iterative, AI-assisted translation process, confirming that a successful migration always necessitates an architectural realignment.